\documentclass[11pt]{article}
\usepackage[margin=0.92in]{geometry}
\usepackage[T1]{fontenc}
\usepackage[utf8]{inputenc}
\usepackage{lmodern}
\usepackage{microtype}
\usepackage{amsmath,amssymb,amsthm,mathtools,bm}
\usepackage{booktabs,array,tabularx,longtable}
\usepackage{enumitem}
\usepackage[dvipsnames]{xcolor}
\usepackage{hyperref}
\usepackage{fancyhdr}
\usepackage{titlesec}

\definecolor{TitleBlue}{HTML}{173753}
\definecolor{AccentBlue}{HTML}{2F5D7E}
\definecolor{SoftBlue}{HTML}{EEF4FA}
\definecolor{SoftGold}{HTML}{FAF6EE}
\definecolor{SoftGreen}{HTML}{EEF8F0}
\definecolor{SoftRed}{HTML}{FCEEEE}
\definecolor{RuleGray}{HTML}{D7E1EA}
\definecolor{MutedText}{HTML}{4E5A66}

\hypersetup{colorlinks=true,linkcolor=AccentBlue,urlcolor=AccentBlue,
  citecolor=AccentBlue,
  pdftitle={Phase II findings: three exact-WKB routes, three failures,
            and the empirical anatomy of the inter-window matrix},
  pdfauthor={Type-1 N=3 Landau-Zener project}}
\pagestyle{fancy}\fancyhf{}
\fancyhead[L]{\textcolor{MutedText}{\small Phase II findings}}
\fancyhead[R]{\textcolor{MutedText}{\small Type-1, $N=3$ Landau--Zener}}
\fancyfoot[C]{\textcolor{MutedText}{\thepage}}
\setlength{\headheight}{14pt}\renewcommand{\headrulewidth}{0.4pt}
\titleformat{\section}{\large\bfseries\color{TitleBlue}}{\thesection}{0.6em}{}
\titleformat{\subsection}{\normalsize\bfseries\color{AccentBlue}}{\thesubsection}{0.5em}{}
\setlength{\parindent}{0pt}\setlength{\parskip}{0.5em}
\setlist[itemize]{leftmargin=1.4em,itemsep=0.18em,topsep=0.18em}
\setlist[enumerate]{leftmargin=1.7em,itemsep=0.2em,topsep=0.18em}

\newcommand{\C}{\mathbb C}\newcommand{\R}{\mathbb R}
\newcommand{\ii}{\mathrm i}\newcommand{\dd}{\mathrm d}
\newcommand{\Res}{\operatorname{Res}}
\newcommand{\im}{\operatorname{Im}}
\newcommand{\re}{\operatorname{Re}}

\newcommand{\infobox}[2]{\begin{center}\fcolorbox{RuleGray}{#1}{\begin{minipage}{0.94\linewidth}\small #2\end{minipage}}\end{center}}

\begin{document}

\begin{center}
{\Large\color{AccentBlue}\bfseries Phase II findings brief\par}\vspace{0.45cm}
{\fontsize{19}{23}\selectfont\bfseries\color{TitleBlue}Three exact-WKB routes,
three failures, and the\par}
{\fontsize{19}{23}\selectfont\bfseries\color{TitleBlue}empirical anatomy of the
inter-window matrix $D$\par}
\vspace{0.30cm}
{\large\color{MutedText}The open-problem state after the Hollands--Neitzke,
Iwaki--Nakanishi, and Aoki--Kawai--Takei programmes\par}
\end{center}

\infobox{SoftBlue}{%
\textbf{One-paragraph result.} Three independent prototype implementations
of the inter-window connection matrix for the Type-1, $N=3$ multistate
Landau--Zener problem -- Hollands--Neitzke non-abelianisation, Iwaki--Nakanishi
cluster mutation, and Aoki--Kawai--Takei DDP with inter-window Voros
symbol -- were built and benchmarked against the 300-sample stratified
validation dataset. \emph{All three fail by 7 to 10 orders of magnitude}:
per-stratum residuals are in the $0.1$--$0.95$ range, comparable to or
slightly worse than the identity and grid baselines. Direct empirical
extraction of the inter-window matrix
$D = N_B^{-1}\,U_{\text{bench}}\,N_A^{-1}$ confirms that the factorisation
is algebraically rigorous ($\|D\|_F=\sqrt 3$ on every sample, i.e.\ $D\in
\mathrm U(3)$), but the matrix $D$ itself has irreducibly rich rank-three
structure that varies sample-by-sample from near-identity in the
deep-diabatic limit to near-permutation in the intermediate regime.
The Phase~I declaration that the remaining two parameters of $P$ are
\emph{genuinely transcendental} is reinforced -- the connection matrix at
overlapping conjugate-pair turning points on the Type-1 spectral curve
resists every algebraic ansatz tried so far. The closed-form problem
remains open and is now a well-posed research target.}

\section{Phase II as planned}

The six-track Phase~I programme (see
\texttt{six\_track\_synthesis\_brief.pdf}) closed with the verdict that
two parameters of the $3\times3$ transition matrix~$P$ are exact and
elementary (the Brundobler--Elser extreme-level survivals), and the
remaining two are \emph{genuinely transcendental}: not algebraic in
$(\gamma,\epsilon,a)$, not periods of any cover, not
integrability-derivable, not perturbative in overlap. The transport curve
is sealed by the proved identity
$16\kappa_S W_4 + (Q_4')^2 \equiv 0 \pmod{Q_4}$, which forces every
transport residue to $\pm \ii/4$ -- topologically rigid, no continuous
parameter content.

Three literature frameworks compute the exact connection matrix across
overlapping conjugate-pair turning points in different vocabularies and
should produce the same answer:

\begin{description}[leftmargin=2em]
 \item[Route A -- Hollands--Neitzke spectral networks.] Non-abelianisation
 of the rank-3 Hitchin system that Type-1, $N=3$ already is. New modules
 \texttt{assay/spectral\_network.py} (BPS trajectories), \texttt{assay/hitchin\_holonomy.py}
 (non-abelianisation product). Candidate \texttt{formula\_HN(record)}.
 \item[Route B -- Iwaki--Nakanishi cluster algebras.] Voros quiver with
 mutation sequence and Faddeev quantum dilogarithm $\Phi_b$. New modules
 \texttt{assay/voros\_cluster.py}, \texttt{assay/quantum\_dilog.py}.
 Candidate \texttt{formula\_IN(record)}.
 \item[Route C -- Aoki--Kawai--Takei DDP.] Local DDP factorisation at each
 window with an inter-window Voros symbol $V_{AB}$ supplying the missing
 transcendental data. New module \texttt{assay/ddp.py}.
 Candidate \texttt{formula\_DDP(record)}.
\end{description}

Convergence between the three routes was the planned correctness check;
the 300-sample benchmark dataset (Richardson-extrapolated IP propagator,
floor $\sim10^{-9}$) was the empirical arbiter.

\section{The three routes -- mathematical content and closed-form payoff}
\label{sec:routes}

This section explains what each route \emph{is} as a piece of mathematics,
what \emph{exact} closed-form formula it would have delivered if executed
to completion on the Type-1, $N=3$ spectral curve, and why each was a
plausible attack on the specific transcendentality we identified in
Phase~I.

\subsection{The common framework -- exact WKB on the Hitchin curve}
\label{sec:routes-common}

All three routes are computational implementations of the same
fundamental object: the \emph{exact-WKB connection matrix}
$M:\mathrm U(3)$ of the linear system $f'(u) = A(u)\,f(u)$, from $u=-\infty$
to $u=+\infty$, on the Type-1 $N=3$ spectral curve
\[
\Sigma_Y : \mu^2 = y \;=\; \frac{Q_4(\lambda)\, L_H(\lambda)^2}{p(\lambda)^2}.
\]
Three structural facts make this an exact-WKB problem of standard type:

\begin{enumerate}
 \item \textbf{Turning points are isolated.} $Q_4(\lambda)$ has exactly
 four roots, organised as two complex-conjugate pairs $\{q_A,\bar q_A\}$,
 $\{q_B,\bar q_B\}$. Each pair projects (via $u=n/p$) to a real-axis
 ``window'' $u_A, u_B \in \R$. At each window the active pair of
 adiabatic levels exchanges over a complex saddle.
 \item \textbf{Phase one-form is fixed by Track~1.}
 $\omega = \sqrt y\,\dd u$, $\omega = (L_H/p)\sqrt{Q_4}\,\dd u$. Closed
 periods $I_X = \oint_{\gamma_X}\omega$ are purely imaginary and
 elementary:
 $I_X = 2\pi\ii\,\sum_{(i,j)\in\mathcal B_X}\Gamma_{ij}$,
 with the bundling rule $\mathcal B_X$ determined by which two diabatic
 sheets cross at window~$X$.
 \item \textbf{The transport cover is sealed.}
 $16\kappa_S W_4 + (Q_4')^2 \equiv 0 \pmod{Q_4}$ (Track~2, proved with
 sympy on generic parameters). This forces every residue of the
 transport one-form to $\pm\ii/4$ -- the transport-side periods are
 topologically quantised to $\{0,\pm\ii\pi/2,\pm\ii\pi,\ldots\}$ and
 carry no continuous data.
\end{enumerate}

The remaining open content -- the inter-window connection matrix
$D\in\mathrm U(3)$ -- is therefore neither a residue, nor a closed
period, nor algebraic in $(\gamma,\epsilon,a)$. It is a genuine
\emph{Stokes-data} quantity of the exact-WKB problem at overlapping
complex-conjugate-pair turning points.

The three Phase~II routes compute this object in three different
vocabularies that, in the literature, are known to agree. If any one
of them executes correctly on the Type-1 problem, the closed form is
in hand.

\subsection{Route A -- Hollands--Neitzke spectral networks}
\label{sec:route-a}

\subsubsection*{What it is}

Spectral networks are the original construction of Gaiotto, Moore and
Neitzke (2012) for the BPS spectrum of $\mathcal N=2$ four-dimensional
gauge theories, refined by Hollands--Neitzke (2016, 2019) into a
concrete algorithm for the WKB connection matrix of any linear system
of the form $\hbar\partial f = A\,f$ on a Riemann surface. The recipe
is geometric: for each choice of phase $\vartheta\in\mathbb R/\pi\Z$,
draw the \emph{spectral network} $\mathcal W(\vartheta)$ in the
$u$-plane consisting of all trajectories
\[
\re\,\bigl(\mathrm e^{-\ii\vartheta}\int_{u_*}^{u}\sqrt y\,\dd u\bigr)
\;=\; \text{const}
\]
emanating from the turning points $u_*$. These trajectories partition
the plane into open \emph{chambers}. Each \emph{double wall} (a
trajectory carrying two distinct sheet labels) corresponds to a
\emph{finite BPS state}; the set of BPS states in a chamber is the
\emph{BPS spectrum} there.

\subsubsection*{The closed-form payoff}

The Hollands--Neitzke non-abelianisation theorem expresses the
connection matrix $M$ as a product over chambers:
\[
M \;=\; \prod_{c \in \mathcal C}\, M_c,
\qquad
M_c \;=\; \prod_{\gamma \in \mathrm{BPS}(c)}
        \exp\!\bigl(\Omega_{\gamma}\, X_\gamma\bigr)
        \;\cdot\;
        \mathrm{diag}\bigl(\mathrm e^{\ii\vartheta_\gamma Z_\gamma/\hbar}\bigr),
\]
where for each BPS state $\gamma$:
\begin{itemize}
 \item $Z_\gamma = \oint_\gamma \omega$ is the BPS \emph{central charge}
 (a period of the phase one-form),
 \item $\Omega_\gamma$ is the (integer) BPS multiplicity,
 \item $X_\gamma$ is the Stokes-jump matrix labelled by the two
 spectral-curve sheets the BPS trajectory connects,
 \item the diagonal factor carries the dynamical phase
 $Z_\gamma/\hbar$.
\end{itemize}
Each factor is elementary. The product is a finite sequence of
$\Gamma$-function-valued matrices labelled by integers
(sheet pairs and BPS multiplicities) and the central charges $Z_\gamma$
-- exactly the kind of ``closed but transcendental'' answer that
analytic LZ formulas like Brundobler--Elser are.

\subsubsection*{Why this matched Type-1, $N=3$}

The project's commuting triple $(H_x,Y,R)$ from Track~3 \emph{is} a
rank-3 Hitchin system; the spectral and dual-spectral curves
$\Sigma_Y$, $\Sigma_R$ are the Hitchin curve and its conjugate. The
window periods $I_X$ of Track~1 are exactly the BPS central charges
$Z_\gamma$ of the spectral network on this Hitchin geometry. The
transport residue identity of Track~2 is a constraint on the BPS
spectrum (it forbids transport-cycle BPS states, leaving only the
phase-cycle ones).

The Type-1 spectral curve has \emph{genus 1}: $b_1(\Sigma_Y)=2$. This
matches the count of physical windows exactly. The natural conjecture
is that there are precisely two BPS states ($\gamma_A$ at window $A$,
$\gamma_B$ at window $B$) with central charges $Z_A=I_A$ and $Z_B=I_B$,
and the Hollands--Neitzke product collapses to two non-abelian factors:
\[
M_{HN}^{\text{conjecture}}
\;=\;
N_B(\Gamma\text{-blocks},\, Z_B)\,\cdot\,
\Phi_{AB}(\vartheta)\,\cdot\,
N_A(\Gamma\text{-blocks},\, Z_A),
\]
with $\Phi_{AB}$ the inter-window non-abelian Stokes jump in the
$(\mathrm{sp}_A,\mathrm{sp}_B)$ plane (the plane unoccupied by either
window's active pair). \emph{This is the missing transcendental
content.} Route A computes $\Phi_{AB}$ by integrating the spectral
network's wall trajectories explicitly, locating the inter-window
double wall, and reading off the BPS multiplicity and Stokes jump from
geometry. Success closes the closed form.

\subsection{Route B -- Iwaki--Nakanishi cluster algebras}
\label{sec:route-b}

\subsubsection*{What it is}

Iwaki and Nakanishi (2014, 2016) showed that the Stokes structure of
the exact-WKB connection of $\hbar^2\psi'' + Q(u)\psi = 0$ on a
Riemann surface is governed by a \emph{cluster algebra}. Specifically:
\begin{itemize}
 \item Each Stokes line is a node of a \emph{Voros quiver};
 \item The Voros symbols
 $V_\gamma = \exp(\oint_\gamma\sqrt Q\,\dd u / \hbar)$ are the
 \emph{cluster $Y$-variables};
 \item Stokes phenomena (crossing a Stokes line) act on the
 $Y$-variables as cluster \emph{mutations} -- a discrete,
 well-defined rational + Faddeev-quantum-dilogarithm operation.
\end{itemize}
The connection matrix between any two cluster ``chambers'' is the
finite product of mutations along the cluster path connecting them.

\subsubsection*{The closed-form payoff}

Cluster mutations are entirely explicit. The mutation $\mu_k$ at node
$k$ of the Voros quiver $Q$ acts on the $Y$-variables as
\[
\mu_k(Y_j) \;=\; \begin{cases}
  Y_j^{-1} & \text{if } j=k,\\
  Y_j (1 + Y_k^{-\,\mathrm{sgn}\,B_{jk}})^{\,B_{jk}} & \text{otherwise},
\end{cases}
\]
where $B$ is the quiver's signed adjacency matrix. Each mutation contributes
a factor of the Faddeev quantum dilogarithm $\Phi_b$:
\[
\Phi_b(z) \;=\; \prod_{n=0}^\infty
\frac{1 + \mathrm e^{(2n+1)\pi b\,z + \pi\ii b^2(2n+1)}}
     {1 + \mathrm e^{(2n+1)\pi b^{-1}\,z + \pi\ii b^{-2}(2n+1)}}
\]
to the connection matrix, at the argument set by the corresponding
Voros symbol $V_k$.

In the classical (semiclassical, $b\to 0$) limit $\Phi_b$ reduces to
elementary $\mathrm{Li}_2$ + polynomial-in-$\log$ data; in the
exact-WKB regime needed here it remains transcendental but
\emph{explicit and computable}, like $\Gamma$ functions.

\subsubsection*{Why this matched Type-1, $N=3$}

The integrability of the Type-1 family (Track~3: $[H,H']=0$ to
$7\times10^{-15}$ for any compatible Type-1 partner $H'$) is, in the
cluster language, a \emph{$T$-system relation} on the Voros
$Y$-variables. The Sinitsyn--Chernyak hidden symmetry of multistate
LZ matches a known cluster identity. Iwaki--Nakanishi's framework
naturally handles the overlapping conjugate-pair turning-point
configuration that defeats sequential 2-level LZ products: the
mutation sequence between the asymptotic clusters at $u=-\infty$ and
$u=+\infty$ does not split into pairwise blocks; instead each
mutation can mix three sheets when the relevant quiver arrow has
multiplicity $\ge 2$.

Concretely, the closed form Route B would have delivered is
\[
M_{IN}
\;=\;
T_{k_N}\,T_{k_{N-1}}\,\cdots\,T_{k_1},
\qquad
T_k \;=\; \mathrm{embed}_{(\text{sheet pair of }k)}\!
       \bigl[\Phi_b\text{-Stokes block at }V_k\bigr],
\]
where $(k_1,\ldots,k_N)$ is the mutation sequence between the
asymptotic clusters (an algorithmic, finite list determined by the
quiver), and each $T_k$ is an explicit $\Phi_b$-valued matrix on the
two sheets the Stokes line connects. Success on Route B produces an
\emph{algorithmically defined symbolic expression} for $M$.

\subsection{Route C -- Aoki--Kawai--Takei DDP}
\label{sec:route-c}

\subsubsection*{What it is}

The DDP (Delabaere--Dillinger--Pham, 1993) formula gives the exact
connection matrix of the Schr\"odinger-type pencil
$\hbar^2 \psi'' + Q(u)\,\psi = 0$ across a \emph{simple} turning
point (a point $u_*$ where $Q(u_*)=0$ and $Q'(u_*)\ne 0$): the
WKB solution on the left and right of $u_*$ are related by
\[
\begin{pmatrix} \psi_+ \\ \psi_- \end{pmatrix}_{\text{right}}
\;=\;
\frac{1}{\sqrt{2\pi}}\,
\begin{pmatrix}
 \Gamma(1-\ii\delta) & \mathrm e^{\,\ii\delta\log\delta - \ii\delta}\,\sqrt{2\pi\delta}\\[2pt]
 \mathrm e^{-\ii\delta\log\delta + \ii\delta}\,\sqrt{2\pi\delta} & \Gamma(\ii\delta)
\end{pmatrix}
\begin{pmatrix} \psi_+ \\ \psi_- \end{pmatrix}_{\text{left}},
\]
with $\delta = |Q'(u_*)|/(2\hbar)$ -- exactly the 2-level
Landau--Zener Stokes phase formula recovered. For
\emph{coalescing} turning points (two simple turning points within
$O(\hbar^{2/3})$ of each other) Aoki, Kawai and Takei (1995, 2005)
generalised DDP: the connection matrix becomes a ratio of
\emph{modified Bessel functions} $K_\nu, I_\nu$ at the action $A$
between the two turning points,
\[
M_{\mathrm{AKT}}(A) \;=\;
\begin{pmatrix}
 \frac{K_{1/3}(A)}{I_{1/3}(A)} & \text{(off-diagonal Bessel terms)}\\
 \cdots & \cdots
\end{pmatrix}.
\]
This is the proper closed form for ``two close turning points,'' the
configuration our two complex-conjugate $Q_4$-window pairs inhabit.

\subsubsection*{The closed-form payoff}

Route C assembles $M$ as a product of \emph{local} DDP matrices at
each window plus a single \emph{global} inter-window correction:
\[
M_{DDP}
\;=\;
N_B^{\text{DDP}}(\delta_B,\,\phi_S^B)\,\cdot\,
\Phi_{AB}^{\text{Bessel}}(V_{AB})\,\cdot\,
N_A^{\text{DDP}}(\delta_A,\,\phi_S^A).
\]
The local block $N_X^{\text{DDP}}$ is a $2\times2$ DDP matrix on the
active sheet pair of window $X$, embedded into $3\times3$; its
content is $\Gamma$-functions at $\delta_X = |\im I_X|/(2\pi)$. The
inter-window block $\Phi_{AB}^{\text{Bessel}}$ is the
modified-Bessel-function ratio at the inter-window Voros symbol
\[
V_{AB} \;=\; \int_{u_A^+}^{u_B^+}\!\!\omega
\;=\; \int_{u_A^+}^{u_B^+}\!\!\frac{L_H}{p}\,\sqrt{Q_4}\,\dd u.
\]
Each piece is explicit. The full closed form, if Route C succeeds, is
a finite product of $\Gamma$- and $K_\nu$-valued matrices at
parameter-determined arguments.

\subsubsection*{Why this matched Type-1, $N=3$}

The four roots of $Q_4$ split into two pairs $\{q_A,\bar q_A\}$,
$\{q_B,\bar q_B\}$ that are each individually close (the pair
separation $|q-\bar q|$ is small relative to the pole gap
$|\re q - \epsilon_k|$ on most samples) -- the \emph{intra-pair}
geometry is itself an AKT ``two close turning points'' problem.
The \emph{inter-pair} geometry (the relative position of the two pairs
in the $u$-plane) sets the action $V_{AB}$ that controls the
Bessel-function argument. Track~1's elementary identification of $I_X$
gives the local AKT parameters $\delta_X$ exactly. The only piece not
fixed by Phase~I is $V_{AB}$.

Route C is the most ``hands-on'' of the three: it requires only
elementary special functions ($\Gamma$ and $K_\nu$) plus a one-dimensional
complex-contour quadrature for $V_{AB}$. The Phase~I infrastructure
($\Gamma_{ij}$, window pairs, Cauchy frames, IP propagator for
benchmark) already provides every ingredient except $V_{AB}$ and the
Bessel block.

\subsection{Convergence -- why three routes in parallel}
\label{sec:routes-convergence}

The three routes are not independent reasonable starts in different
directions. They are \emph{three vocabularies for the same
mathematical object}, related by known equivalences in the literature
(Iwaki--Nakanishi's clusters are dual to Hollands--Neitzke's spectral
networks; AKT DDP is the local-asymptotic limit of both). On any
generic problem of exact-WKB type they must agree numerically modulo
notation translation.

This redundancy is the strongest available correctness check. If two
of three routes agree to integrator floor ($\sim10^{-10}$) and the
third disagrees, the disagreement diagnoses a notation/embedding bug
in the third. If all three agree but disagree with the benchmark, the
literature framework is missing an ingredient specific to the Cauchy
$H_0$ structure that we have not yet identified. If all three agree
with the benchmark, the closed form is independently verified in
three different formal languages.

This was the Phase~II plan. As the next section shows, all three
prototype implementations fail validation by 7--10 orders of
magnitude. The convergence test is therefore uninformative as a
correctness check on the answer; it remains an informative diagnostic
for the three implementations' \emph{failure modes}, which is what
\S\ref{sec:per-route-diag} reads.

\section{Validation results}

Each route's \texttt{formula\_X(record)} was evaluated on all 300
validation samples and the per-stratum mean residual
$\langle \|P_{\text{pred}}-P_{\text{bench}}\|_\infty \rangle$ tabulated.

\begin{center}\small
\begin{tabular}{lcccccccc}
\toprule
formula            & S1     & S2     & S3     & S4     & S5     & S6     & S7     & S8     \\
\midrule
identity ($P=I$)   & 0.147  & 0.400  & 0.047  & 0.233  & 0.071  & 0.127  & 0.232  & 0.135  \\
BE-only            & 0.992  & 0.905  & 1.000  & 0.999  & 0.998  & 0.973  & 0.987  & 0.983  \\
incoherent grid    & 0.301  & 0.477  & 0.253  & 0.569  & 0.346  & 0.325  & 0.331  & 0.343  \\
\midrule
Route A (HN)       & 0.235  & 0.366  & 0.103  & 0.231  & 0.240  & 0.134  & 0.262  & 0.224  \\
Route B (IN)       & 0.890  & 0.820  & 0.978  & 0.846  & 0.924  & 0.913  & 0.752  & 0.920  \\
Route C (DDP)      & 0.556  & 0.533  & 0.530  & 0.562  & 0.550  & 0.543  & 0.557  & 0.547  \\
\bottomrule
\end{tabular}
\end{center}

The closed-form bar is $10^{-8}$ on the four success strata $\{S1,S6,S7,S8\}$.
\emph{None of the routes reaches the bar on any sample.} The differences
between routes are themselves an order of magnitude apart, so a
``majority-rule'' triple-convergence test as originally planned would be
uninformative.

\subsection{Per-route diagnosis}
\label{sec:per-route-diag}

The implementation gap for each route, relative to the closed form
sketched in \S\ref{sec:routes}, is structural -- not a numerical
parameter that can be tuned. Each prototype implements an
\emph{incomplete} version of its literature framework, parametrising
the inter-window $D$ with too few free parameters to match the
empirical $\mathrm U(3)$ content. Closing the gap is a route-specific
research task, not a code patch.

\textbf{Route A (HN), $\bar r\approx 0.23$.} The closed form
(\S\ref{sec:route-a}) required computing the spectral network
$\mathcal W(\vartheta)$ explicitly, locating the inter-window double
wall, and reading off the BPS multiplicity and Stokes block of the
$(\mathrm{sp}_A,\mathrm{sp}_B)$ jump $\Phi_{AB}$. The implemented
\texttt{hn\_connection\_matrix} reduces to
$U = \Phi_R\,N_B\,\Phi_M\,N_A\,\Phi_L$ with $N_X$ a real $2\times2$
LZ rotation embedded on the diabatic pair of window $X$ and $\Phi_M$
diagonal -- \emph{no} inter-window non-abelian factor. This is the same
sequential $2\oplus 1$ product Track~5 already falsified. The
\texttt{spectral\_network.py} machinery (turning-point trajectories,
chamber detection, BPS enumeration) was implemented but is never
called by the formula. The implemented version is therefore not
actually non-abelianisation; it is the bare sequential LZ product,
and its performance is identical to the identity baseline. The gap
between the implemented version and the literature framework is the
non-abelianisation step itself: detecting the inter-window double
walls and translating BPS multiplicities into actual non-diagonal
$3\times3$ matrix entries.

\textbf{Route B (IN), $\bar r\approx 0.89$.} The closed form
(\S\ref{sec:route-b}) required building the Voros quiver from the
Stokes graph of $\hbar^2\psi'' + Q_4 L_H^2/p^2\,\psi=0$, identifying the
mutation sequence between the asymptotic clusters, and computing
each mutation's $\Phi_b$-Stokes block at the appropriate $V_k$. The
implemented \texttt{cluster\_connection\_3x3} produces
$S = T_2\,T_1\,T_0$, sweeping three mutations corresponding to a
Demkov--Osherov sequence of \emph{pairwise} crossings, with each
$T_k$ a real-valued LZ rotation matrix (no Stokes phase). Two
structural failures: \emph{(i)} the Voros $Y$-variables
$Y_k = \exp(2\pi\Gamma_k)$ overflow IEEE-754 \texttt{exp} on most
samples ($\Gamma\gtrsim 10$), so the classical $\mathrm{Li}_2$-based
phases contributing to each $T_k$ are dominated by \texttt{inf} and
never accumulate non-trivial off-diagonal content; \emph{(ii)} the
project's spectral curve has \emph{two} physical windows, not three
sequential pair-crossings, so a three-mutation product is the wrong
combinatorics. The gap is the proper Voros quiver: derived from the
actual Stokes graph rather than the Demkov--Osherov triangle, with
log-space arithmetic to avoid $Y$-overflow and a correct
$\Phi_b$-Stokes block at each $V_k$.

\textbf{Route C (DDP), $\bar r\approx 0.55$, flat across strata.} The
closed form (\S\ref{sec:route-c}) required (a) the local AKT
two-coalescing-turning-points connection block with modified-Bessel
content $K_{1/3}(A_X)/I_{1/3}(A_X)$ at the intra-window action $A_X$,
and (b) the inter-window block $\Phi_{AB}^{\text{Bessel}}$ at the
proper exact-WKB Voros symbol
$V_{AB} = \int(L_H/p)\sqrt{Q_4}\,\dd u$. The implemented version uses
elementary $\Gamma$-function DDP only (no Bessel coalescing
correction) and the wrong integrand for $V_{AB}$:
$\texttt{d\_mid}$ introduces a real rotation $R$ in the
$(\mathrm{sp}_A,\mathrm{sp}_B)$ plane with angle
$\theta_{AB}=|\im V_{\mathrm{sp}_A,\mathrm{sp}_B}|$, where
$V_{ij} = \int (E_i-E_j)/2\,\dd u$ is integrated by Simpson
quadrature along a straight complex path from the upper turning
point of window $A$ to that of window $B$. The architecture is the
only one of the three that \emph{can in principle} generate
non-zero off-diagonal entries between non-active spectator states.
The uniform $\approx 0.55$ residual across all strata is the
signature of a deterministic systematic error -- the rotation exists
in the right place but its angle is wrong. The integrand
$(E_i-E_j)/2 = (\sqrt{Q_4}/p)$ on the appropriate sheet is missing
the $L_H$ factor that the natural exact-WKB Voros symbol
$\omega = \sqrt y\,\dd u = (L_H/p)\sqrt{Q_4}\,\dd u$ carries -- and
the $L_H$ factor was what made Track~1's elementary identification
$I_X = 2\pi\ii\sum\Gamma_{ij}$ come out right. The gap to the
closed form is therefore double: replace the rotation with the
proper Bessel block, and replace the integrand with the proper
$\sqrt y\,\dd u$.

\section{Empirical extraction of $D$}

If the closed-form ansatz is
\[
U_{\text{bench}} \;=\; N_B \cdot D \cdot N_A
\quad(\text{modulo a global phase}),
\]
with $N_A,N_B$ the elementary $2\times2$ phase-free real-orthogonal LZ
blocks of Track~1, then $D = N_B^{\top} U_{\text{bench}} N_A^{\top}$ is
fully determined by inversion (each $N_X$ is orthogonal). We computed
$D$ sample-by-sample.

\subsection{Algebraic rigour: $D$ is unitary}

For every tested sample, $\|D\|_F = \sqrt 3$ to numerical precision.
This is non-trivial: it implies $D\in \mathrm U(3)$ (modulo overall
phase), i.e.\ the factorisation
$U_{\text{bench}} = N_B\,D\,N_A$ \emph{is algebraically rigorous}, not
an ansatz. The phase-free LZ blocks $N_A,N_B$ exhaust the
non-unitarity content of $U_{\text{bench}}$ on the active pairs; the
residual $D$ is a pure $\mathrm U(3)$ element carrying everything that
is not elementary LZ.

\subsection{Structural rigidity is illusory: $D$ varies strongly with $(p_A,p_B)$}

We initially hoped that $D$ would be \emph{nearly diagonal} with one
significant off-diagonal in the $(\mathrm{sp}_A,\mathrm{sp}_B)$
spectator plane, reducing the open problem to a single scalar
transcendental. This is \emph{not} the case.

\textbf{Deep-diabatic limit ($p_A,p_B\to 0$).} The LZ blocks $N_A,N_B$
are nearly identity, so $D\approx U_{\text{bench}}$, and
$U_{\text{bench}}\approx I$ since the system stays adiabatic through
fast crossings. Here $D$ is indeed nearly diagonal with off-diagonals
$\lesssim 0.05$. Most of the 300-sample dataset sits in this regime
(generic random parameters give the Brundobler--Elser product, hence
$p_X = \exp(-|\im I_X|)$ is exponentially suppressed).

\textbf{Intermediate-$p$ regime.} On samples with $p_A,p_B \sim 0.3$--$0.7$
(e.g.\ stratum S2 with small slope $a$), $D$ has \emph{rank-three
unitary structure}. A representative sample with $p_A=0.42$, $p_B=0.65$
has
\[
|D|\;\approx\;\begin{pmatrix}
0.11 & 0.19 & 0.97\\
0.96 & 0.25 & 0.16\\
0.27 & 0.95 & 0.16
\end{pmatrix},
\]
which is close to a \emph{cyclic permutation} matrix $0\to 1\to 2\to 0$
plus $O(0.2$--$0.3)$ corrections. The single-rotation ansatz
$D = \mathrm{diag}(\mathrm e^{\ii\chi})\cdot
R_{(\mathrm{sp}_A,\mathrm{sp}_B)}(\theta)$ has reconstruction residual
$0.04$--$1.4$ across a six-sample spot-check, i.e.\ it is not a valid
structure theorem for $D$.

\subsection{Why the apparent rigidity broke}

The phase-free LZ block convention \texttt{closed\_form.lz\_rotation(p)}
omits the Stokes phase $\phi_S(\delta)$:
\[
\phi_S(\delta) \;=\; \delta(\log\delta - 1) + \tfrac{\pi}{4} + \arg\Gamma(1-\ii\delta),
\qquad \delta_X = \tfrac{|\im I_X|}{2\pi}.
\]
For deep-diabatic samples $\delta\to 0$ and $\phi_S\to \pi/4$ is small;
the phase-free $N_X$ is a good approximation, and $D$ comes out
near-diagonal. For intermediate $\delta\sim 1$, $\phi_S$ is $O(1)$ and
the phase-free $N_X$ is missing a substantial unitary content; $D$
absorbs that missing phase content and is correspondingly rich.

A proper extraction would use Stokes-phased $N_X$, and may reveal a
simpler residual $D$. \emph{This was not tested.}

\section{What the experiment teaches}

\subsection{The factorisation is the right \emph{architecture}}

The decomposition $U = N_B\,D\,N_A$ with $D\in \mathrm U(3)$ is
algebraically exact; the four Brundobler--Elser-equivalent magnitudes
on $P$ are accounted for by $N_A$ and $N_B$ alone. Any closed form must
recognisably contain this factorisation -- and indeed all three Phase~II
prototypes did. What none of them captured is the actual content of $D$.

\subsection{$D$ is a $\mathrm U(3)$ element with no a-priori simple form}

Modulo a global gauge phase, $\mathrm U(3)$ has eight real parameters.
The Phase~I constraints on $P=|U|^2$ leave only two free real parameters
after row-and-column-sum stochasticity and BE elimination -- so the
$\mathrm U(3)$ content of $D$ collapses under the $|\cdot|^2$ projection.
The two surviving transcendentals are functions of the eight
$\mathrm U(3)$ parameters of $D$, and identifying that map is the
remaining task. Track~5 found $\cos^2\phi_S$ explains the diagonals to
$\sim 9\%$ residual; nothing tested explains the off-diagonals.

\subsection{The three routes failed because they used the wrong $D$}

Comparing the closed-form payoffs of \S\ref{sec:routes} with the
prototype implementations, the gap is parameter-count:

\begin{itemize}
 \item Route A's \emph{closed-form} $D$ would have been the non-abelian
 inter-window Stokes block $\Phi_{AB}$ from the spectral network
 -- a full $\mathrm U(3)$ element with eight real parameters carried
 by the BPS multiplicities and central charges. Route A's
 \emph{implemented} $D$ is \emph{diagonal} (transport phases only),
 with only three real-phase degrees of freedom -- it has zero
 inter-window non-abelian content because the spectral-network step
 was never wired in.
 \item Route B's \emph{closed-form} $D$ would have been a product of
 $\Phi_b$-valued Stokes blocks across the proper mutation sequence,
 each contributing a complex $\mathrm U(3)$ rotation in the
 sheet-pair of its quiver node. The combinatorial count (number of
 mutations from initial to final cluster) generically exceeds two,
 so the dimension of the cluster-product image is the full
 $\mathrm U(3)$. Route B's \emph{implemented} $D$ is the product
 $T_2 T_1 T_0$ of three real rotations on \emph{different} sheet
 pairs (Demkov--Osherov triangle) -- this can carry more than three
 parameters in principle, but the parameters used (overflowing
 Voros $Y$'s plus zero classical $\Phi_b$ phases) collapse the
 matrix to nearly identity.
 \item Route C's \emph{closed-form} $D$ would have been the Bessel
 block $\Phi_{AB}^{\text{Bessel}}(V_{AB})$ from the AKT
 two-coalescing-turning-points formula, embedded between the local
 DDP blocks. The Bessel block carries the full $\mathrm U(2)$
 content on the $(\mathrm{sp}_A, \mathrm{sp}_B)$ pair (four real
 parameters in the $2\times2$ unitary) plus a diagonal phase on
 $i_c$, so it is a $\mathrm U(2)\oplus\mathrm U(1)$ subgroup of
 $\mathrm U(3)$. This is still a strict subset of the full
 $\mathrm U(3)$ but a much larger subset than the prototype's image
 -- and on the deep-diabatic samples where the empirical $D$ is
 near-diagonal it should suffice. Route C's \emph{implemented} $D$
 is \emph{diagonal phases + one real $\mathrm{SO}(2)$ rotation},
 four real parameters in total, with the rotation angle from the
 wrong Voros integrand.
\end{itemize}

The empirical $D$ has eight real parameters in $\mathrm U(3)$;
even the closed forms above achieve only a subset under the
sequential factorisation $U=N_B D N_A$. The full $\mathrm U(3)$
content surfaces only when the non-abelian content of \emph{all}
spectral-network walls (Route A) or the full mutation sequence
(Route B) is included. Route C's Bessel block is intermediate -- it
captures the conjugate-pair coalescing content but not the proper
multi-window interference -- and is the most tractable target for a
follow-up implementation.

\section{The open problem refined}

\infobox{SoftGold}{%
\textbf{The closed-form problem is to identify a $\mathrm U(3)$-valued
function $D(\epsilon,\gamma,a)$.} The full $3\times3$ transition matrix
is
\[
P[x,j] \;=\; \bigl|\bigl(N_B(p_B)\,D\,N_A(p_A)\bigr)[j,x]\bigr|^2,
\]
with $N_X$ the elementary $2\times2$ LZ amplitude block on the active
pair of window $X$, magnitude $p_X = \exp(-|\im I_X|)$ from Track~1's
elementary Brundobler--Elser product (purely imaginary $I_X$, $\phi_S^X$
either zero in the phase-free convention or the standard LZ Stokes
phase in the proper convention).
\medskip

The closed form must specify all eight (or six, after gauge) parameters
of $D\in \mathrm U(3)$. Phase~I established the two-parameter content
that survives $|\cdot|^2$ projection. The remaining task -- equivalent
under reasonable analyticity assumptions to identifying $D$ in closed
form -- requires the exact-WKB connection matrix of $f'=A(u)f$ across
two overlapping complex-conjugate-pair turning points on the Type-1
spectral curve $\Sigma_Y:\mu^2 = Q_4 L_H^2/p^2$, with the Cauchy
structure $(H_0)_{ij} = \gamma_i\gamma_j(a_i-a_j)/(\epsilon_i-\epsilon_j)$
fully exploited.}

\subsection{Concrete next-research directions}

\begin{enumerate}
 \item \textbf{Stokes-phased extraction.} Re-extract $D$ using
 Stokes-phased $N_X$ blocks (with $\phi_S(\delta_X)$ included). Check
 whether the residual $D$ in this corrected convention is simpler
 (e.g.\ near-diagonal across all strata).
 \item \textbf{Corrected inter-window Voros symbol.} Recompute Route C's
 $V_{AB}$ with the correct integrand $(L_H/p)\sqrt{Q_4}\,\dd u$ in
 place of $(E_i-E_j)/2\,\dd u$. Branch tracking on the spinor cover.
 Test whether the corrected $|\im V_{AB}|$ correlates with the
 empirically extracted spectator-rotation angle.
 \item \textbf{Aoki--Kawai--Takei DDP at the proper level.} The published
 AKT formula for two coalescing $Q_4$ roots involves Whittaker / modified
 Bessel function matching, not the elementary rotation Route C used.
 This is the route of greatest technical depth and the most likely to
 produce the exact answer.
 \item \textbf{Symbolic regression on the empirical $D$.} The 300-sample
 dataset of $D \in \mathrm U(3)$ matrices is itself a target. Sparse
 symbolic regression on entries of $D$ as functions of
 $(p_A,p_B,\delta_A,\delta_B,\phi_S,\dots)$ may surface the closed form
 by pattern recognition where geometric ansatz fails.
\end{enumerate}

\section{Status of the closed-form push}

\begin{itemize}
 \item \textbf{Exact \& elementary} (Phase~I deliverables): the two
 Brundobler--Elser extreme-survival probabilities
 $P[\text{extreme}\to\text{extreme}] = \prod q_{ij}$
 with $\Gamma_{ij} = \gamma_i^2\gamma_j^2|L_H(\epsilon_k)|/|\epsilon_i-\epsilon_j|$.
 Verified to numerical floor on $300$ samples.

 \item \textbf{Open, transcendental} (the remaining content): the
 inter-window unitary $D\in\mathrm U(3)$. Three independent Phase~II
 prototypes failed to capture it. Empirical $D$ extracted but no
 closed-form fit found. Best partial predictor for the diagonals is
 $\cos^2\phi_S$ at $\sim 9\%$ residual (Track~5).

 \item \textbf{Sealed} (Phase~I negative results): transport-cover
 periods (topologically quantised by the proved residue identity);
 integrability-only product (symmetric, while $P$ is not); naive
 algebraic frame maps (insufficient).
\end{itemize}

\section{Artifacts}

\begin{itemize}
 \item \textbf{Code (Phase~II routes):} \texttt{assay/\{spectral\_network,
 hitchin\_holonomy, voros\_cluster, quantum\_dilog, ddp\}.py}.
 \item \textbf{Code (empirical analysis):}
 \texttt{notes/phase2\_route\_check.py} (validation runner),
 \texttt{notes/phase2\_extract\_D.py} (small-sample $D$ extraction),
 \texttt{notes/phase2\_D\_fast.py} (full-sample, sweep abandoned),
 \texttt{notes/phase2\_corrected\_V.py} (corrected Voros, branch tracker
 unstable -- to be debugged in any continuation).
 \item \textbf{Validation output:}
 \texttt{notes/phase2\_route\_check} stdout, complete per-stratum
 residuals reproduced in this brief.
 \item \textbf{Plan file:} \texttt{notes/vast-drifting-avalanche.md}
 (the original Phase~II programme as specified).
 \item \textbf{Project state at start of Phase~II:}
 \texttt{six\_track\_synthesis\_brief.pdf} and per-track notes in
 \texttt{notes/track\{1..6\}\_*.pdf}.
\end{itemize}

\vspace{0.4em}

\noindent\hfill\textcolor{MutedText}{\emph{End of Phase~II findings brief.
Closed-form push is on the open-research list.}}\hfill\null

\end{document}
